% sections/03_science.tex -- Q1/Q2/Q3 科学分析

\section{科学分析：太空散热与热环境}

\subsection{Q1：真空辐射散热的物理极限}

\subsubsection{斯特藩-玻尔兹曼定律基础}

在真空中，热辐射是唯一的散热途径。斯特藩-玻尔兹曼定律决定了单位面积散热通量的上限：

\begin{equation}
\frac{P}{A} = \varepsilon \sigma T^{4}
\label{eq:sb}
\end{equation}

其中 $\sigma = 5.6704\times10^{-8}$\unitWpm{}/K\textsuperscript{4}（CODATA 2018），$\varepsilon$ 为表面辐射率（0--1），$T$ 为散热面绝对温度（K）。

\begin{figure}[H]
\centering
% figures/sb_t4_curve.tex
% 斯特藩-玻尔兹曼 T⁴ 曲线：散热通量 vs 温度
\begin{tikzpicture}
\begin{axis}[
    width=0.95\textwidth,
    height=7cm,
    xlabel={\textbf{散热面温度} $T$ (K)},
    ylabel={\textbf{散热通量} (W/m\textsuperscript{2})},
    xmin=290, xmax=510,
    ymin=0, ymax=3800,
    grid=major,
    legend style={
        at={(0.02,0.98)},
        anchor=north west,
        font=\footnotesize,
        cells={anchor=west}
    },
    tick label style={font=\footnotesize},
    label style={font=\small},
]

% ε=0.85
\addplot[color=gray!60, line width=1.2pt, mark=*, mark size=1.5pt]
    coordinates {(300,390)(323,523)(340,642)(350,723)(373,932)(400,1234)(450,1981)(500,3015)};
\addlegendentry{$\varepsilon=0.85$}

% ε=0.90
\addplot[color=accent, line width=1.2pt, mark=square*, mark size=1.5pt]
    coordinates {(300,413)(323,555)(340,682)(350,766)(373,988)(400,1306)(450,2093)(500,3190)};
\addlegendentry{$\varepsilon=0.90$}

% ε=0.95
\addplot[color=primary!70, line width=1.2pt, mark=triangle*, mark size=1.5pt]
    coordinates {(300,436)(323,586)(340,720)(350,808)(373,1043)(400,1379)(450,2209)(500,3367)};
\addlegendentry{$\varepsilon=0.95$}

% ε=1.00 (黑体)
\addplot[color=primary, line width=1.5pt, dashed, mark=diamond*, mark size=1.5pt]
    coordinates {(300,459)(323,617)(340,758)(350,851)(373,1097)(400,1452)(450,2325)(500,3544)};
\addlegendentry{$\varepsilon=1.00$ (黑体)}

% AI1 工作点标注
\addplot[only marks, mark=*, mark size=4pt, color=highlight]
    coordinates {(342.5,1400)};
\node[highlight, anchor=south west, font=\footnotesize\bfseries]
    at (axis cs:342.5,1400) {AI1 工作点 (342.5\,K, $\sim$1400\,W/m\textsuperscript{2})};

% T⁴ 非线性辅助标注
\draw[->, >=stealth, graybase, line width=0.6pt]
    (axis cs:380,800) -- (axis cs:460,1800);
\node[graybase, font=\footnotesize, anchor=north]
    at (axis cs:420,1050) {\textbf{$T^{4}$ 非线性增长}};

\end{axis}
\end{tikzpicture}

\caption{斯特藩-玻尔兹曼 $T^{4}$ 曲线：散热通量随温度的非线性增长。4条曲线对应不同辐射率（$\varepsilon=0.85/0.90/0.95/1.00$），红色标记为AI1工作点（342.5\,K，$\sim$1400\,W/m\textsuperscript{2}）。注意$T^{4}$效应使高温区通量急剧上升。}
\label{fig:sb_t4}
\end{figure}

表\ref{tab:sb_basic}给出不同温度和辐射率下的单面散热通量：

\begin{table}[H]
\centering
\footnotesize
\caption{斯特藩-玻尔兹曼定律基础散热通量（单面辐射，W/m\textsuperscript{2}）}
\label{tab:sb_basic}
\begin{tabular}{rrrrrr}
\toprule
$T$ (K) & $T$ (\unitC{}) & $\varepsilon=0.85$ & $\varepsilon=0.90$ & $\varepsilon=0.95$ & $\varepsilon=1.00$ \\
\midrule
300 & 27 & 390 & 413 & 436 & 459 \\
323 & 50 & 523 & 554 & 585 & 616 \\
340 & 67 & 642 & 680 & 718 & 756 \\
350 & 77 & 723 & 765 & 808 & 851 \\
373 & 100 & 932 & 987 & 1,042 & 1,097 \\
400 & 127 & 1,234 & 1,306 & 1,379 & 1,452 \\
450 & 177 & 1,981 & 2,098 & 2,214 & 2,331 \\
500 & 227 & 3,015 & 3,192 & 3,370 & 3,547 \\
\bottomrule
\end{tabular}
\end{table}

\subsubsection{双面辐射与AI1 1,400 W/m\textsuperscript{2} 验证}

AI1采用双面辐射散热板设计——物理面板面积110 m\textsuperscript{2}，等效辐射面积220 m\textsuperscript{2}（两面辐射至深空）。双面辐射通量（按物理面板面积计）为：

\begin{equation}
\frac{P}{A_{\text{物理面板}}} = 2\varepsilon\sigma T^{4}
\label{eq:sb_double}
\end{equation}

关键计算结果（表\ref{tab:ai1_verify}）：

\begin{table}[H]
\centering
\caption{AI1散热参数反推与物理可行性验证}
\label{tab:ai1_verify}
\begin{tabular}{cccc}
\toprule
\textbf{辐射率 $\varepsilon$} & \textbf{双面模式所需 $T$ (K)} & \textbf{所需 $T$ (\unitC{})} & \textbf{验证通量 (W/m\textsuperscript{2})} \\
\midrule
0.85 & 347.9 & 74.8 & 1,400 \\
0.88 & 345.0 & 71.9 & 1,400 \\
\textbf{0.90} & \textbf{342.5} & \textbf{69.5} & \textbf{1,403} \\
0.92 & 340.2 & 67.1 & 1,400 \\
0.95 & 336.8 & 63.7 & 1,400 \\
\bottomrule
\end{tabular}
\end{table}

\textbf{AI1 1,400 W/m\textsuperscript{2} 物理可行性判定：物理成立，已接近当前材料天花板。}

在双面辐射、$\varepsilon\approx0.90$、面板温度约342.5K（69.5\unitC{}）时可达到1,400\unitWpm{}（物理面板面积）。面板温度69.5\unitC{}在液冷散热器工作范围内。但$\varepsilon=0.90$需要高性能热控涂层（如Z-93白漆），在空间环境退化（UV+原子氧）后的EOL（寿命末期）$\varepsilon$会下降至约0.82--0.85，届时散热能力将衰减5--10\%\cite{nasa2024sst,nasa1993thermal}。

\subsubsection{1 GW轨道数据中心所需散热面积}

假设1 GW中$\sim$95\%转化为废热（950 MW待散热），关键场景的散热面积需求见表\ref{tab:onegw_area}：

\begin{table}[H]
\centering
\footnotesize
\caption{1 GW轨道数据中心所需最小散热面积（多场景）}
\label{tab:onegw_area}
\begin{tabular}{lcccrr}
\toprule
\textbf{场景} & $\varepsilon$ & $T$ (K) & \textbf{模式} & \textbf{通量 (W/m\textsuperscript{2})} & \textbf{面积 (km\textsuperscript{2})} \\
\midrule
AI1 BOL（激进） & 0.90 & 342.5 & 双面 & 1,403 & 0.71 \\
AI1 EOL（退化后） & 0.83 & 340 & 双面 & $\sim$1,250 & $\sim$0.80 \\
典型航天器设计点 & 0.85 & 350 & 双面 & 1,446 & 0.69 \\
保守设计点 & 0.85 & 323 & 单面 & 523 & 1.91 \\
高性能涂层+高温 & 0.93 & 353 & 双面 & 1,740 & 0.57 \\
极限材料+最高温 & 0.96 & 393 & 双面 & 2,970 & 0.34 \\
\bottomrule
\end{tabular}
\end{table}

\textbf{核心洞察}：AI1设计点的1 GW散热面积约0.71 km\textsuperscript{2}（$\sim$700,000 m\textsuperscript{2}，相当于边长840 m正方形）。按每颗AI1卫星110 m\textsuperscript{2}散热板估算，需要约6,400颗AI1级卫星来散热1 GW\cite{nasa2024sst,cybernative2026thermo}。

\subsection{Q2：太空vs地面散热效率定量对比}

\subsubsection{PUE对比：太空的实际优势}

PUE（Power Usage Effectiveness）衡量数据中心冷却开销。表\ref{tab:pue_compare}汇总了地面与太空的PUE对比：

\begin{table}[H]
\centering
\caption{太空vs地面数据中心PUE与冷却效率对比}
\label{tab:pue_compare}
\begin{tabular}{lccc}
\toprule
\textbf{来源} & \textbf{PUE} & \textbf{冷却开销（\%总功率）} & \textbf{年份} \\
\midrule
行业平均（Uptime Institute） & 1.56--1.58 & 35.9--36.7 & 2024 \\
Google全球机队平均 & 1.09 & 8.3 & 2024 \\
最佳超大规模 & 1.06--1.12 & 5.7--10.7 & 2024 \\
GB300液冷配置（理论） & 1.05--1.10 & 4.8--9.1 & 2025--26 \\
\rowcolor{gray!10}
\textbf{太空辐射散热（理论）} & \textbf{$\sim$1.02--1.05} & \textbf{2--5（仅泵功率）} & \textbf{--} \\
\bottomrule
\end{tabular}
\end{table}

\textbf{太空的核心PUE优势不是"散热更快"，而是"不需要冷却塔/制冷机组等额外能耗设施"——冷源（2.7K宇宙微波背景）天然存在且免费。}\cite{uptime2024,google2024pue,microsoft2024pue}

\subsubsection{散热密度对比：太空的天然劣势}

然而，辐射散热通量（单面$\sim$400--1,500\unitWpm{}）比地面液冷冷板（10,000--100,000\unitWpm{}）低1--2个数量级。辐射散热受限于$T^{4}$效应——即使将面板温度从350K提升至500K（+43\%），通量也仅从$\sim$800提升至$\sim$3,000\unitWpm{}（+275\%），但500K（227\unitC{}）已超出绝大多数电子设备的安全工作范围。

\subsubsection{黄仁勋"97.5\%重量是冷却系统"说法的量化验证}

多源验证表明：GB300 NVL72机架总重约1.36--1.59吨（而非2吨）；机架内液冷组件约占7--12\%（$\sim$100--170 kg）；加上CDU（冷却分配单元）后冷却相关总重量约18--35\%\cite{lenovo2025gb300,aivres2025krs}。黄仁勋的说法是\textbf{修辞性夸张，方向正确但数字夸大}——97.5\%更可能指传统风冷数据中心中冷却基础设施占总设施重量的高比例（含冷却塔、冷水机组等），而非单个机架的冷却组件。

\subsection{Q3：轨道热环境的真实复杂性}

\begin{figure}[H]
\centering
% figures/leo_thermal_env.tex
% LEO 轨道热环境示意图 —— 纯 TikZ
\begin{tikzpicture}[
    >=Stealth,
    scale=1.0,
    every node/.style={font=\footnotesize},
]

% ===== 中心卫星 =====
\filldraw[fill=primary!10, draw=primary, line width=1.5pt]
    (-1.5,-0.8) rectangle (1.5,0.8);
\node[primary, font=\bfseries] at (0,0) {\textbf{AI1 散热板}};
\node[primary, font=\scriptsize] at (0,-0.5) {(辐射散热至深空 2.7\,K)};

% ===== 太阳直射（上方金黄色，从太阳指向卫星） =====
\draw[->, line width=1.5pt, color=yellow!80!orange]
    (0,3.5) -- (0,0.8);
\node[anchor=south, text=yellow!60!orange, font=\footnotesize\bfseries]
    at (0,3.6) {\textbf{太阳直射}};
\node[anchor=north, text=yellow!70!orange, font=\scriptsize]
    at (0,2.8) {1361\,W/m\textsuperscript{2}};
\node[anchor=north, text=yellow!70!orange, font=\scriptsize]
    at (0,2.3) {吸收$\sim$5\% (knife-edge 定向后)};

% ===== 地球红外辐射（下方蓝色，从地球指向卫星） =====
\draw[->, line width=1.2pt, color=blue!60]
    (-1.2,-3.0) -- (-0.3,-0.8);
\node[anchor=north, text=blue!60, font=\footnotesize\bfseries]
    at (-2.0,-3.0) {\textbf{地球红外}};
\node[anchor=south, text=blue!60, font=\scriptsize]
    at (-1.2,-2.3) {$\sim$240\,W/m\textsuperscript{2}};

% ===== 地球反照（下方浅蓝，从地球指向卫星） =====
\draw[->, line width=1.2pt, color=cyan!60]
    (1.2,-3.0) -- (0.3,-0.8);
\node[anchor=north, text=cyan!60, font=\footnotesize\bfseries]
    at (2.0,-3.0) {\textbf{地球反照}};
\node[anchor=south, text=cyan!60, font=\scriptsize]
    at (1.2,-2.3) {$\sim$480\,W/m\textsuperscript{2}};

% ===== 辐射散热至深空（左右灰色箭头） =====
\draw[->, line width=1.5pt, color=graybase]
    (-1.5,0) -- (-3.5,0);
\node[anchor=east, text=graybase, font=\footnotesize\bfseries]
    at (-3.6,0.3) {\textbf{辐射散热}};
\node[anchor=east, text=graybase, font=\scriptsize]
    at (-3.6,-0.25) {1400\,W/m\textsuperscript{2}};

\draw[->, line width=1.5pt, color=graybase]
    (1.5,0) -- (3.5,0);
\node[anchor=west, text=graybase, font=\footnotesize\bfseries]
    at (3.6,0.3) {\textbf{辐射散热}};
\node[anchor=west, text=graybase, font=\scriptsize]
    at (3.6,-0.25) {1400\,W/m\textsuperscript{2}};

% ===== 深空标注 =====
\node[graybase, font=\scriptsize\itshape]
    at (3.8,2.5) {深空背景};
\node[graybase, font=\scriptsize\itshape]
    at (3.8,2.1) {2.7\,K};

% ===== 地球示意 =====
\filldraw[fill=blue!20, draw=blue!40, line width=0.8pt]
    (-3.0,-4.0) arc (180:360:3.0 and 0.8) -- (3.0,-4.0) -- cycle;
\node[blue!60, font=\scriptsize\bfseries] at (0,-4.0) {\textbf{地球}};

\end{tikzpicture}

\caption{LEO轨道热环境示意图。中央为AI1散热板（双面辐射至2.7\,K深空），四周展示四种热流：太阳直射（1361\,W/m\textsuperscript{2}，knife-edge定向后吸收$\sim$5\%）、地球红外辐射（$\sim$240\,W/m\textsuperscript{2}）、地球反照（$\sim$480\,W/m\textsuperscript{2}）以及向深空的辐射散热（1400\,W/m\textsuperscript{2}，双面）。}
\label{fig:leo_thermal_env}
\end{figure}

\subsubsection{LEO热环境全景参数}

LEO并非"理想冷库"，而是高度动态的多热源环境。表\ref{tab:leo_thermal}汇总了关键热环境参数\cite{clawson2002thermal,nasa2021tm4527}：

\begin{table}[H]
\centering
\footnotesize
\caption{LEO轨道关键热环境参数}
\label{tab:leo_thermal}
\begin{tabular}{lcc}
\toprule
\textbf{参数} & \textbf{值} & \textbf{波动范围} \\
\midrule
太阳常数（1 AU均值） & 1,367\unitWpm{} & 1,322--1,414（$\pm$3.5\%） \\
地球反照率均值 & $\sim$0.30 & 0.05（海面）--0.85（雪/云） \\
反照热流（LEO典型） & $\sim$480\unitWpm{}（局部最大值） & 视地表类型大幅波动 \\
地球红外辐射（OLR）均值 & $\sim$240\unitWpm{} & 220--260 \\
轨道周期 & $\sim$90分钟 & 85--95（400--800 km） \\
日照/阴影比 & $\sim$60/40 & 取决于$\beta$角 \\
外部表面温度范围 & -150\unitC{}至+150\unitC{} & 峰谷差可达220\unitC{} \\
\bottomrule
\end{tabular}
\end{table}

\subsubsection{散热器效率衰减量化}

在多热源真实环境中，散热器净效率显著低于理想真空计算值\cite{modus2026thermal,encyclopedia2025leo}：

\begin{table}[H]
\centering
\footnotesize
\caption{LEO环境对散热器效率的衰减因素}
\label{tab:efficiency_loss}
\begin{tabular}{p{2.8cm}p{5.2cm}p{1.8cm}p{2.8cm}}
\toprule
\textbf{衰减因素} & \textbf{衰减机制} & \textbf{效率损失} & \textbf{应对策略} \\
\midrule
太阳直射吸收 & 散热面朝向太阳时吸收$\sim$1,367\unitWpm{} & 100\%（完全丧失） & knife-edge定向+太阳盾 \\
\hline
太阳掠射（knife-edge, $<$5\unitC{}入射角） & 极小投影面积$\times$太阳常数 & $\sim$5--10\%净散热损失 & 主动姿态控制$\pm$1\unitC{} \\
\hline
地球反照加热 & 反照流$\sim$480\unitWpm{}$\times$VF$\times\alpha_s$ & $\sim$10--20\% & 散热面朝向深空 \\
\hline
地球红外加热 & OLR$\sim$240\unitWpm{}$\times$VF$\times\varepsilon_{IR}$ & $\sim$5--10\% & 同上 \\
\hline
$\beta$角高值期间 & 轨道日照比例高，散热器长时间暴露于太阳 & $\sim$15--30\%（时段性） & 择机部署 \\
\hline
热控涂层退化（EOL） & UV+原子氧导致$\alpha_s$上升、$\varepsilon_{IR}$下降 & $\varepsilon$下降5--15\% & 选择抗退化涂层 \\
\bottomrule
\end{tabular}
\end{table}

\subsubsection{knife-edge-to-sun定向有效性}

knife-edge策略（散热器法线矢量与太阳方向夹角=90\unitC{}）可将太阳吸收降至散热能力的$<$1\%（姿态控制精度$\pm$1\unitC{}时，太阳吸收功率$\approx24\times\alpha_s\approx3.6$\unitWpm{}，vs散热功率$\sim$800\unitWpm{}）。但要求持续主动姿态控制，且$\beta$角不利期间效率显著降低。

\textbf{LEO昼夜循环影响}：90分钟周期（日照$\sim$54分钟/阴影$\sim$36分钟）对百kW级载荷不可通过PCM（相变材料）储热解决——缓冲150kW$\times$36分钟需$\sim$1.9吨纯PCM\cite{ijisrt2025pcm}。必须依赖散热器自然热容+液冷环路热容平滑+适度降载策略。

\textbf{综合结论}：散热器效率在真实LEO环境中比理想真空计算值低15--40\%，具体取决于轨道$\beta$角、散热面定向策略和涂层退化程度。

% end of 03_science
